\begin{figure}[H]
\centering
\begin{minipage}[t]{0.4\textwidth}
\vspace{0pt}\centering
\includegraphics[width=\linewidth]{figures/ch01/panels/fig1_1a_rotation.pdf}
\end{minipage}\hfill
\begin{minipage}[t]{0.4\textwidth}
\vspace{0pt}\centering
\includegraphics[width=\linewidth]{figures/ch01/panels/fig1_1b_shear.pdf}
\end{minipage}
\caption{\textbf{基方向决定线性变换。}
浅色虚线表示原始坐标网格，实线表示变换后的网格。箭头 $A\hat e_1$ 与 $A\hat e_2$ 是变换后的基方向，也就是矩阵 $A$ 的两列。旋转保持方形网格，而剪切会使一族网格线发生倾斜；一旦变换后的基方向确定，其他任意向量的变换都可由线性组合得到。}
\label{fig:ch1-matrix-transformations}
\end{figure}